% Supplement subsections typed below as "Section~S<n>.<m>"; build.sh checks each number against the built supplement.
% supplement-ref supp:learning=S3.5
% supplement-ref supp:borda=S1.8
\section{Related Work}\label{sec:related}

ArrowFlow touches seven lines of work. The first studies ordinal codes and distances between rankings, and the second studies voting rules that combine rankings. The third covers prototype methods, and the fourth covers kernels, metric learning and symbolic vectors. The fifth covers alternatives to backpropagation, and the sixth covers ways to train through a sort. The seventh supplies the tools that ArrowFlow's pipeline borrows. For each line we say what it does and where ArrowFlow departs from it.

\paragraph{Ordinal codes and distances between rankings.}
The idea that order alone can carry information is old. Rank-order coding holds that the order in which neurons fire carries information about the stimulus \citep{thorpe1998rapid,gautrais1998rate,thorpe2001spike}, and \citet{bonilla2022analyzing} compare it with codes based on the time of the first spike. The rank transform of nonparametric statistics \citep{conover1981rank} also keeps only order. But it ranks one variable across observations, whereas top-scoring-pair classifiers compare pairs of features within one example \citep{geman2004classifying}. Rank-ordered autoencoders order hidden activations and learn sparse codes through progressive reconstruction \citep{bertens2016rank}. ArrowFlow instead feeds an ordinal code into layers of learned ranking filters and compares rankings by Spearman's footrule. \citet{diaconis1977footrule} show that the footrule stays within a factor of two of Kendall's distance, the number of item pairs that two rankings order differently \citep{kendall1938}. \citet{fagin2003topk} and \citet{fagin2006partial} extend both distances to top-$k$ lists and to partial rankings. Learning to rank \citep{burges2005learning,burges2010ranknet} fits a real-valued scoring function whose output is an ordering. In ArrowFlow, the inputs and the parameters are themselves orderings.

\paragraph{Voting rules and rank aggregation.}
Combining many rankings into one is the subject of social choice theory. Arrow's theorem \citep{arrow1951social,arrow1963social} limits every rule that turns a profile of rankings, one per voter, into one ranking. Positional rules, such as the Borda count \citep{young1974borda}, give each item points for its position on each ballot. Distance rules, such as the Kemeny median \citep{kemeny1959mathematics} and the footrule median, pick the ranking with the smallest total distance to the ballots. Both kinds have been compared in rank aggregation for metasearch \citep{dwork2001rank}. The Mallows model \citep{mallows1957nonnull} is a probability distribution over rankings that is centered on one modal ranking. \citet{conitzer2005voting} show which voting rules are maximum-likelihood estimators under which noise models, so choosing a rule also chooses what it estimates. Under Mallows noise, positional rules such as the Borda count recover the modal ranking as the number of votes grows \citep{caragiannis2013noisy}. The Borda count also combines the class rankings of several classifiers \citep{ho1994decision,vanerp2000variants}. ArrowFlow applies a weighted Borda count inside a learning system. Section~\ref{sec:arrow} states what Arrow's theorem implies for the network and what it does not.

\paragraph{Prototypes and reference objects.}
Many classifiers predict from distances to learned reference patterns. Learning vector quantization (LVQ) and self-organizing maps work this way \citep{kohonen1990self}, and so do generalized LVQ \citep{sato1995glvq,hammer2014learning,nebel2015median} and adaptive resonance theory \citep{grossberg1987competitive}. \citet{diao2021glvq} combine random-feature encodings with generalized LVQ and connect the result to hyperdimensional classifiers. ArrowFlow's output layer, which drives training, is a prototype classifier of this kind, with permutations as prototypes. Describing an object by its distances to reference objects is also the dissimilarity representation of \citet{pekalska2005dissimilarity}. Metric-space search goes one step further. It represents an object by the order of the reference objects by their distance to it, and it compares such orders with rank distances \citep{chavez2008proximity,amato2008approximate}. An untrained ranking layer read by a footrule nearest-neighbor vote is such a permutation index, with random rankings instead of data points as references. ArrowFlow learns the reference permutations and stacks the construction in layers. Random-weight networks fix random hidden units and train only a readout. Random vector functional-link networks \citep{pao1994learning} and extreme learning machines \citep{huang2006extreme,huang2012extreme} are examples, and random features approximate kernel machines in the same spirit \citep{rahimi2007random}. An untrained ArrowFlow is a permutation analog of these models.

\paragraph{Kernels, metric learning and symbolic vectors.}
Other methods also classify permutation-valued inputs, but they do not learn the permutations themselves. Kernels on permutations are the established supervised route. \citet{jiao2015kernels,jiao2018kendall} classify with Kendall and Mallows kernels, including numerical data converted to rankings, and they learn a decision function from pairwise similarities. Supervised metric learning fits a representation for a nearest-neighbor classifier \citep{goldberger2004nca,weinberger2009lmnn}. ArrowFlow instead trains its permutations against one prototype per class. Vector symbolic architectures and hyperdimensional computing (HDC) represent structured objects in high-dimensional distributed vectors. Holographic reduced representations bind and superpose vectors with circular convolution \citep{plate1995hrr}. Kanerva's construction labels an item's position before summation so that order survives \citep{kanerva2009hdc}, and ordered position codes serve biosignal classification \citep{rahimi2018biosignal}. The surveys of \citet{kleyko2022survey1,kleyko2023survey2} map the field. In these systems the state is distributed and algebraic. The codes are still related to ArrowFlow's. A thermometer code writes position $p$ as $p$ ones followed by zeros, and the footrule equals the Hamming distance between the thermometer codes of the positions. But ArrowFlow stores explicit permutations and updates them by position votes. Learned hyperdimensional classifiers train their class vectors \citep{duan2022lehdc,smets2025margin}, and \citet{verges2025hdcclassification} review and compare such classifiers. So the hyperdimensional control of Section~S3.5, an untrained classifier built from fixed random codes of items and positions, represents only its own construction. It says nothing about learned hyperdimensional classifiers. Discrete encoders can be trained with native integer and binary operations, with reported gains over untrained ones \citep{kirby2026encoder}.

\paragraph{Alternatives to backpropagation.}
The forward-forward algorithm \citep{hinton2022forward}, target propagation \citep{bengio2014auto} and its difference form \citep{lee2015difference}, and feedback alignment \citep{lillicrap2016random} avoid sending exact gradients back through a network, but they keep real-valued weights. ArrowFlow replaces the weights by permutations and updates them with position votes. Each batch update exactly minimizes a squared position objective of one filter (Section~S1.8), but no loss of the whole network is minimized. The output layer's vote only attracts the filter of the true class, whereas LVQ1 and LVQ2.1 also push a prototype of a wrong class away \citep{kohonen1990self}. A hidden layer computes the signal for the layer below from the permutations of its own selected filters. This is the analog of weight transport, the reuse of the forward weights to send errors back, which feedback alignment avoids. We do not propose the rule as a model of learning in the brain.

\paragraph{Training through a sort.}
Differentiable sorting replaces the permutation with a smooth approximation so that gradients can pass through it. NeuralSort \citep{grover2019neuralsort}, SoftSort \citep{prillo2020softsort}, projection onto the permutahedron \citep{blondel2020fast}, optimal transport \citep{cuturi2019ot} and Gumbel-Sinkhorn networks \citep{mena2018gumbelsinkhorn}, built on Sinkhorn normalization \citep{sinkhorn1967diagonal}, all work this way. ArrowFlow keeps the sort exact and never differentiates through it. Its trained encoder (Section~\ref{sec:learned-encoder}) also learns through an exact sort, by turning the motions into target scores in the spirit of target propagation \citep{bengio2014auto,lee2015difference}. Other methods also keep a discrete step exact. The straight-through estimator copies a gradient across the step as if the step were the identity \citep{bengio2013estimating,sahoo2023identity}. Blackbox differentiation interpolates the loss behind the solver, using a second call to the solver \citep{vlastelica2020differentiation,rolinek2020optimizing}. Perturbed optimizers average the solver's output over random perturbations of its input \citep{berthet2020perturbed}, and implicit maximum-likelihood estimation builds a target distribution for it \citep{niepert2021imle}. LambdaRank sets score gradients from desired rank changes \citep{burges2006nonsmooth}. The structured perceptron and its generalization, the Fenchel-Young losses, train a scoring model on the exact output of a solver \citep{collins2002perceptron,blondel2020fenchel}. Target propagation for hard-threshold units sets discrete targets for the hidden units and trains each unit toward them \citep{bengio2014auto,friesen2018combinatorial}. The trained encoder's targets come from motions toward a learned class filter and away from the nearest wrong one, and the example's own sorted scores turn them back into scores.

\paragraph{Random projections and ensembles.}
ArrowFlow's fixed encoder rests on random projection. The known guarantees on preserved distances \citep{johnson1984extensions,achlioptas2003database} and preserved margins \citep{arriaga2006algorithmic} concern the projected vector before sorting. The order of two Gaussian-projected coordinates is a random-hyperplane bit \citep{charikar2002similarity}, the sign of a random linear function of the input. The order of Gaussian projections also underlies the rank-order hashes for cosine similarity of \citet{eshghi2008locality}, and it relates to winner-take-all hashing \citep{yagnik2011power}. Section~\ref{sec:theory} states the resulting angle property. The multi-view ensemble is related to random-projection ensemble classification \citep{cannings2017random}, although ArrowFlow neither selects its projections nor thresholds the vote. It stands in the tradition of bagging \citep{breiman1996bagging}, boosting \citep{schapire1990strength}, ensemble diversity \citep{dietterich2000ensemble,kuncheva2003diversity} and multi-view learning \citep{xu2013survey}.
